\section{Билет 56. Распределение Максвелла молекул по скоростям и его экспериментальное подтверждение (опыт Ламмерта). Распределение молекул по энергии.}

\begin{definition}
Экспериментальное подтверждение распределения Максвелла было получено в опытах О. Ламмерта (1920--1930-е гг.). Установка состояла из двух быстро вращающихся соосных цилиндров со щелями, смещёнными на угол $\varphi$. Молекулы из источника проходили через первую щель и достигали детектора только если их скорость удовлетворяла условию:
\begin{equation}
    v = \frac{l \omega}{\varphi},
\end{equation}
где $l$ --- расстояние между цилиндрами, $\omega$ --- угловая скорость вращения. Изменяя $\omega$ или $\varphi$, можно выделять молекулы с разными скоростями. Измеренное распределение интенсивности молекулярного пучка полностью совпало с теоретическим предсказанием Максвелла, подтвердив правильность функции $F(v) \propto v^2 e^{-mv^2/(2kT)}$.
\end{definition}

\begin{theorem}[Распределение по кинетической энергии]
Перейдём от распределения по скоростям к распределению по кинетической энергии поступательного движения $\varepsilon = \frac{1}{2} m v^2$. Тогда $v = \sqrt{2\varepsilon/m}$ и $dv = \frac{d\varepsilon}{\sqrt{2m\varepsilon}}$. Подставляя в формулу Максвелла $\frac{dN}{N} = F(v) dv$, получаем:
\begin{equation}
    \frac{dN}{N} = f(\varepsilon) \, d\varepsilon = \frac{2}{\sqrt{\pi}} (kT)^{-3/2} \sqrt{\varepsilon} \, \exp\left( -\frac{\varepsilon}{kT} \right) d\varepsilon. \label{eq:energy_dist_56}
\end{equation}
Функция $f(\varepsilon)$ описывает долю молекул, энергия которых лежит в интервале $(\varepsilon, \varepsilon+d\varepsilon)$. Распределение нормировано на единицу.
\end{theorem}

\begin{property}[Наиболее вероятная и средняя энергия]
Наиболее вероятная энергия $\varepsilon_{\text{вер}}$ находится из условия максимума $df/d\varepsilon = 0$:
\begin{equation}
    \varepsilon_{\text{вер}} = \frac{1}{2} kT. \label{eq:most_prob_energy_56}
\end{equation}
\textbf{Важно:} $\varepsilon_{\text{вер}} \neq \frac{1}{2} m v_{\text{вер}}^2 = kT$. Это связано с нелинейным преобразованием переменных и наличием множителя $\sqrt{\varepsilon}$ в распределении.
Средняя кинетическая энергия вычисляется интегрированием:
\begin{equation}
    \langle \varepsilon \rangle = \int_0^\infty \varepsilon f(\varepsilon) \, d\varepsilon = \frac{3}{2} kT. \label{eq:avg_energy_56}
\end{equation}
Результат согласуется с законом равнораспределения энергии по трём поступательным степеням свободы.
\end{property}

\begin{remark}
Распределение по энергии имеет практическое значение в химической кинетике и физике фазовых переходов. Доля молекул с энергией, превышающей некоторый порог $\varepsilon_0$ (например, энергию активации реакции или работу выхода при испарении), определяется интегралом $\int_{\varepsilon_0}^\infty f(\varepsilon) d\varepsilon \propto e^{-\varepsilon_0/(kT)}$. Экспоненциальная зависимость от $1/T$ объясняет сильное ускорение химических реакций и процессов испарения при нагреве (уравнение Аррениуса).
\end{remark}
